% !PRINTED 274
\weberpage{274}

We consider first the so-called \gesperrt{spherical functions}, which are used
in many ways in mathematical physics and mechanics.

By spherical functions one understands a system of integral rational
functions of increasing degree, which is defined as follows:
\[
\begin{array}{l}
P_0(x) = 1 \\
P_1(x) = x \\
P_2(x) = \frac{3}{2}(x^2 - \frac{1}{3}) \\
P_3(x) = \frac{5}{2}(x^3 - \frac{3}{5}x) \\
\hbox{\quad\quad . \quad . \quad . \quad . \quad . \quad . \quad . \quad . \quad .}
\end{array}
\]
\[
P_n(x) = \frac{1.3 \ldots (2n-1)}{1.2 \ldots n}
\left(x^n - \frac{n(n-1)}{2(2n-1)}x^{n-2}
+ \frac{n(n-1)(n-2)(n-3)}{2.4(2n-1)(2n-3)}x^{n-4} - \cdots\right),
\]
so that $P_n(x)$ is an integral rational function of degree $n$,
and, according as $n$ is even or odd, contains only the even or
only the odd powers. By use of the sign $\Pi(n)$ one can also
represent the general expression for the spherical functions thus:
\weq{1}{P_n(x) = \sum^{\mu} \frac{(-1)^\mu}{2^n}
\frac{\Pi(2n-2\mu)x^{n-2\mu}}{\Pi(\mu)\,\Pi(n-2\mu)\,\Pi(n-\mu)},}
in which the sum with respect to $\mu$ is to be extended from
$\mu = 0$ so far as $n - 2\mu$ does not become negative.

Between these functions there hold the following relations, which,
after substituting the expression (1), can be verified by simple
comparison of the coefficients of equal powers of $x$,
\weq{2}{\begin{aligned}
n P_n(x) - (2n-1) x P_{n-1}(x) + (n-1) P_{n-2}(x) &= 0 \\
P_1(x) - x P_0(x) &= 0
\end{aligned}}
\weq{3}{(1-x^2) P_n'(x) + n x P_n(x) - n P_{n-1}(x) = 0.}
From the last equation it follows for $x = \pm 1$
\[
\pm P_n(\pm 1) = P_{n-1}(\pm 1),
\]
hence
\weq{4}{P_n(1) = 1, \quad P_n(-1) = (-1)^n.}
From (2) and (3) it now follows that the functions $P_n$, ordered
downwards from an arbitrary $m$,
%CONT>>
